% ID: SEL-D0002-OA-E03
% ESTADO: cotejada. Fuente y transcripción contrastadas de forma independiente; salida revisada.
% Fuente: Selectividad/sel_2019_matematicas_aplicadas/Titular Septiembre/Exámenes Mat Apl CCSS II - A y B.pdf
% SHA-256: 31fb46947e141f5c8da8f7b381e7d27b67fe5a9d1597120454ce66fce0313f1a
% Página: 1; opción A; ejercicio 3.
% Encabezado original: EJERCICIO 3. Solo enunciado; no contiene solución.

Una marca de patinetes eléctricos fabrica tres modelos distintos A, B y C.
El modelo A supone el 25\,\% de su producción, el B el 40\,\% y el resto
de la producción corresponde al modelo C. Transcurridos tres meses desde
su venta, se comprobó que el 15\,\% de patinetes del modelo A, el 10\,\%
del B y el 12\,\% del C había presentado alguna avería. Se elige al azar
un patinete de esta marca.

\begin{enumerate}[label=\alph*)]
  \item \textbf{(1 punto)} Calcule la probabilidad de que dicho patinete haya
  presentado alguna avería.
  \item \textbf{(0.5 puntos)} Si sabemos que el patinete elegido es del
  modelo A, \textquestiondown{}cuál es la probabilidad de que no haya
  presentado avería?
  \item \textbf{(1 punto)} Calcule la probabilidad de que haya presentado
  avería o sea del modelo C.
\end{enumerate}
